% !TEX root = ../main.tex

\chapter{引言}\label{ch:introduction}
  \section{研究背景与意义}\label{sec:background}
    偏微分方程 (Partial Differential Equations, PDEs) 是处理各种物理, 金融问题的重要手段.
    近年来计算机硬件和计算科学的发展使得低维的简单问题基本可以快速求解.
    但是由于传统的有限元方法 (Finite Element Method, FEM) 在解决高维度问题时网格数量呈指数增长, 时间复杂度随之爆炸,
    出现了人们常说的维度诅咒 (curse of dimensionality) 现象, 因此需要找到更高效的方法求解高维 PDEs.
    蒙特卡洛采样 (Monte Carlo sampling, MC sampling) 通过避免全空间的网格生成解决了时间复杂度问题,
    而前馈神经网络 (Feedforward Neural Networks, FNNs) 则解决了在蒙特卡洛采样下编码边界条件与优化解函数的问题.
    近年来, 机器学习 (Machine Learning, ML) 在高维函数逼近中的表现引起了科学计算和计算数学领域的关注.
    E\cite{eMachineLearningComputational2020} 在相关综述中指出,
    传统有限差分方法 (Finite Difference Method, FDM), FEM 和谱方法 (Spectral Method, SM)
    已经能够常规处理三维空间加时间的问题,
    稀疏网格方法可以进一步将可处理维度提高到约 $8$ 到 $10$ 维,
    但对于上百维 PDE, 传统网格型方法仍然难以直接应用.
    因此, 利用深度学习或其他随机化函数逼近方法求解高维 PDE,
    成为近几年高维数值计算中的一个重要方向.
    本文则在 BSDE 求解框架下使用随机特征法 (Random Feature Method, RFM) 代替常规神经网络,
    以降低参数优化的复杂度.

  \section{抛物型方程的基本类型}\label{sec:ppde}
    抛物型方程是一类重要的偏微分方程, 其典型形式为
    \begin{equation}
      \frac{\partial u}{\partial t} = \mathcal{L} u + f,
      \label{eq:ppde0}
    \end{equation}
    其中 $\mathcal{L}$ 为关于空间变量的二阶线性椭圆算子.
    该类方程通常用于描述随时间演化的扩散过程, 例如热传导, 概率密度演化以及金融衍生品定价等问题.
    为了方便与倒向随机微分方程 (Backward Stochastic Differential Equations, BSDEs) 对齐, 我们将椭圆算子拆分得到如下形式:
    \begin{equation}
      \begin{aligned}
        \frac{\partial u}{\partial t}(t, x)
        &+ \frac{1}{2} \operatorname{Tr}\!\big(
        \sigma\sigma^{\operatorname{T}}(t, x)\, \nabla_x^2 u(t, x) \big) \\
        &+ \langle \mu(t, x), \nabla u(t, x) \rangle \\
        &+ f\!\big( t, x, u(t,x), \sigma^{\operatorname{T}}(t, x) \nabla u(t, x) \big)
        = 0.
      \end{aligned}
      \label{eq:ppde}
    \end{equation}

    \subsection{线性抛物型方程}\label{subsec:lppde}
      首先考虑其中的线性方程, 即 $f$ 是 $u, \nabla u$ 的线性函数.
      一方面, 从数学结构上看, 线性抛物型方程具有良好的正则性与稳定性,
      其解通常可以通过解析表示 (如基本解或 Green 函数) 或概率表示 (如 Feynman--Kac 公式) 来刻画,
      可以作为 baseline 用来判断方法的正确性并评估计算效率.
      另一方面, 线性的方程也可以像半线性方程那样通过 Itô 公式变换为 BSDE, 使得此方法可以简单的推广到非线性情形.
      下面列举了常见的线性抛物型方程:
      \begin{enumerate}
        \item
        在本文的统一半线性抛物型方程框架下, 热方程对应的随机过程为标准 $d$ 维布朗运动,
        即漂移项为零, 扩散系数为单位矩阵. 因此有
        \begin{equation*}
          \mu(t, x) = 0,\qquad
          \sigma(t, x) = I_{d},\qquad
          f(t, x, u, z) = 0.
        \end{equation*}
        代入统一形式后得到
        \begin{equation*}
          \frac{\partial u}{\partial t}(t,x)
          + \frac{1}{2}\Delta u(t,x)
          =0,
          \qquad
          u(T,x)=g(x).
        \end{equation*}
        热方程是最经典的线性抛物型偏微分方程之一, 最初来源于热传导问题,
        也广泛用于描述粒子扩散, 浓度传播以及随机扰动在空间中的扩散过程.
        从概率角度看, 热方程与布朗运动具有直接联系, 其解可以表示为布朗运动终端函数的期望.
        因此, 热方程不仅是扩散模型的基础方程, 也是检验随机方法求解线性抛物型 PDE 的标准基准问题.
        在高维情形下, 热方程的维度对应空间变量的维数, 即 $x\in\mathbb{R}^d$ 中的 $d$.
        该算例主要用于验证随机特征法在纯扩散线性问题上的基本逼近能力.

        \item
        在统一框架下, 多资产 Black--Scholes 方程对应
        \begin{equation*}
          \mu(t, x) = r\,x,\qquad
          \sigma(t, x) = \operatorname{diag}(x)\, \Sigma^{1/2},\qquad
          f(t, x, u, z) = -r\,u.
        \end{equation*}
        代入统一形式后得到
        \begin{equation*}
          \frac{\partial u}{\partial t}(t,x)
          + \frac{1}{2}\operatorname{Tr}\!\big(
            \operatorname{diag}(x)\Sigma\operatorname{diag}(x)\nabla_x^2 u(t,x)
          \big)
          + \langle r x,\nabla_x u(t,x)\rangle
          - r u(t,x)
          =0,
          \qquad
          u(T,x)=g(x).
        \end{equation*}
        Black--Scholes 方程来源于金融衍生品定价理论, 是数学金融中最具有代表性的线性抛物型方程之一.
        Black 和 Scholes\cite{scholesPricingOptionsCorporate} 在期权定价研究中将期权价格视为标的资产价格与时间的函数,
        并通过无套利思想和动态对冲方法推导出相应的定价方程;
        在经典单资产情形下, 该模型给出了欧式期权的封闭形式定价公式,
        即把期权价值写成股票价格和时间的函数, 并通过构造对冲组合得到定价微分方程.
        在多资产情形下, 若期权收益依赖于多个标的资产, 则 PDE 的空间维度等于资产数量,
        因而 Black--Scholes 方程自然成为高维 PDE 方法中的典型应用算例.
        设 $x=(x_{1},\ldots,x_{d})$ 表示 $d$ 个资产价格, $r$ 表示无风险利率, $\Sigma$ 表示资产收益率的协方差矩阵.
        在风险中性测度下, 资产价格通常服从几何布朗运动, 其漂移项为 $r x$, 扩散项与当前资产价格成正比.
        其中 $-r u$ 项表示未来收益在无风险利率下的折现效应.
        该算例用于检验本文方法在具有乘性扩散, 非零漂移和折现项的高维金融定价问题上的表现.

        \item
        Ornstein--Uhlenbeck 方程对应的随机过程是一类典型的均值回复过程,
        最早常用于描述带阻尼的随机运动, 后来也广泛出现在统计物理, 金融数学和随机控制等领域.
        与布朗运动不同, Ornstein--Uhlenbeck 过程的漂移项会将状态变量拉回长期均值附近,
        因而适合描述具有稳定中心或长期均衡水平的随机系统.
        在高维情形下, 其维度同样对应状态变量的空间维数.

        \item
        线性反应--扩散方程用于描述扩散过程与局部反应过程同时存在的现象.
        其中扩散项刻画物质, 能量或信息在空间中的传播, 反应项刻画局部增长, 衰减或外部源项作用.
        这类方程在化学反应, 生物种群动力学, 热源扩散以及物理场演化问题中都有广泛应用.
        在本文中考虑的线性反应--扩散方程仍然采用标准扩散过程作为其随机过程部分, 即漂移项为零, 扩散系数为单位矩阵.
        与热方程不同的是, 方程中额外加入了关于解函数 $u$ 的线性反应项以及可能依赖于时间和空间的外部源项 $c(t,x)$.
      \end{enumerate}

      尽管部分方程在物理建模中通常定义于低维空间 (如三维), 但在金融建模, 多因子随机过程等问题中,
      其状态变量维度可以显著增加, 因此常被用作高维数值算法的基准问题.
      由于前两者已经能包含大部分线性方程的形式, 且后面还有非线性方程的实验,
      本论文将仅对热方程和 Black--Scholes 方程进行数值实验.

    \subsection{半线性抛物型方程}\label{subsec:slppde}
      如果 $f$ 不是 $u$ 与 $\nabla u$ 的线性函数, 则整个方程为半线性抛物型方程.
      尽管半线性方程解的解析表示不那么明显, 但通常也有较为完善的理论对其数值做近似估计.
      在完成线性方程求解的理论分析与实验验证后, 对半线性方程的实验可以进一步说明此方法的有效性.
      作为线性的推广, 我们依然通过 Itô 公式将此类方程变换为 BSDE 进行求解.
      下面列举了常见的半线性抛物型方程:
      \begin{enumerate}
        \item
        本文采用线性二次型高斯控制问题对应的 Hamilton--Jacobi--Bellman 方程
        (Hamilton--Jacobi--Bellman equation, HJB equation).
        在统一半线性抛物型方程框架下, 取
        \begin{equation*}
          \mu(t, x) = 0,\qquad
          \sigma(t, x) = \sqrt{2} I_{d},
        \end{equation*}
        此时
        \begin{equation*}
          z = \sigma^{\operatorname T} \nabla_{x} u = \sqrt{2} \nabla_{x} u.
        \end{equation*}
        若控制强度参数为 $\lambda>0$, 生成元取为
        \begin{equation*}
          f(t, x, y, z) = - \frac{\lambda}{2} \lVert z \rVert^{2},
        \end{equation*}
        则得到具体方程
        \begin{equation*}
          \frac{\partial u}{\partial t}(t, x)
          + \Delta u(t, x)
          - \lambda \lVert \nabla_{x} u(t, x) \rVert^{2}
          = 0,
          \qquad
          u(T, x) = g(x).
        \end{equation*}
        Hamilton--Jacobi--Bellman 方程来源于动态规划和随机最优控制理论.
        在最优控制问题中, 解函数通常表示从当前状态出发所能达到的最优代价, 因而也称为值函数.
        当状态变量维数较高时, HJB 方程的维度等于控制系统的状态空间维度, 传统网格方法会迅速受到维数灾难的限制.
        因此, HJB 方程是检验高维 PDE 方法的重要非线性算例.
        从控制问题角度看, 可设状态变量 $X_{t}\in\mathbb{R}^{d}$ 满足受控随机微分方程
        \begin{equation*}
          \dif X_{t} = 2\sqrt{\lambda}\, m_{t} \dif t + \sqrt{2} \dif W_{t},
        \end{equation*}
        其中 $m_{t}$ 为控制变量, $\lambda > 0$ 表示控制强度, $W_{t}$ 为 $d$ 维标准布朗运动.
        控制目标为最小化代价泛函
        \begin{equation*}
          J(\{m_{t}\}_{0 \leq t \leq T}) =
          \mathbb{E}\left[ \int_{0}^{T} \lVert m_t\rVert^{2} \dif t + g(X_{T})\right].
        \end{equation*}
        由动态规划原理可得到上面的 HJB 方程.
        该算例主要用于检验本文方法处理梯度型非线性项的能力.

        \item
        在统一半线性抛物型方程框架下, Allen--Cahn 方程可取
        \begin{equation*}
          \mu(t, x) = 0,\qquad
          \sigma(t, x) = \sqrt{2} I_{d},
        \end{equation*}
        并令
        \begin{equation*}
          f(t, x, y, z) = y - y^{3}.
        \end{equation*}
        代入统一形式后得到终端值问题
        \begin{equation*}
          \frac{\partial u}{\partial t}(t, x)
          + \Delta u(t, x)
          + u(t, x) - u^{3}(t, x)
          = 0,
          \qquad
          u(T, x) = g(x).
        \end{equation*}
        Allen--Cahn 方程是一类典型的非线性反应--扩散方程, 常用于描述相分离, 界面演化以及有序--无序转变等物理现象.
        其核心特点是同时包含扩散项和双势阱型反应项.
        扩散项刻画空间传播或平滑效应, 非线性反应项刻画系统向不同稳定相的演化趋势.
        与 HJB 方程不同, Allen--Cahn 方程的非线性主要依赖于解函数 $u$ 本身, 而不是梯度.
        对应的正向时间形式可写为
        \begin{equation*}
          \frac{\partial v}{\partial s}(s, x) = \Delta v(s,x) + v(s,x) - v^3(s,x),
          \qquad
          v(0, x) = g(x).
        \end{equation*}
        为了与本文采用的终端值问题形式一致, 令
        \begin{equation*}
          u(t, x) = v(T-t, x),
        \end{equation*}
        对给定初始点 $x_0$, 需要近似的目标量为
        \begin{equation*}
          u(0, x_0),
        \end{equation*}
        它对应原始正向时间方程中 $v(T, x_0)$ 的值.
        该算例主要用于检验本文方法对反应型非线性项的处理能力.

        \item
        在统一框架下, 带违约风险的非线性 Black--Scholes 方程可对应
        \begin{equation*}
          \mu(t, x)= \bar{\mu}x, \qquad
          \sigma(t, x) = \bar{\sigma} \operatorname{diag}(x),
        \end{equation*}
        以及
        \begin{equation*}
          f(t, x, y, z) = - (1 - \delta) Q(y) y - R y.
        \end{equation*}
        代入统一形式后得到
        \begin{equation*}
          \frac{\partial u}{\partial t}(t, x)
          + \bar{\mu} x \cdot \nabla_{x} u(t, x)
          + \frac{\bar{\sigma}^{2}}{2}
          \sum_{i = 1}^{d} x_{i}^{2}
          \frac{\partial^{2} u}{\partial x_{i}^{2}}(t, x)
          - (1 - \delta) Q \big( u(t, x) \big) u(t, x)
          - R u(t, x)
          =0.
        \end{equation*}
        带违约风险的非线性 Black--Scholes 方程是金融衍生品定价中的一个重要扩展模型.
        经典 Black--Scholes 方程假设市场不存在违约风险, 此时方程是线性的.
        但在实际金融市场中, 合约发行方可能发生违约, 合约持有人只能获得部分价值补偿.
        若违约强度依赖于当前合约价值, 则定价方程中会出现关于 $u$ 的非线性项.
        设 $x=(x_{1}, \ldots, x_{d})$ 表示 $d$ 个标的资产价格, 在风险中性测度下资产价格可建模为几何布朗运动.
        其中 $R$ 为无风险利率, $\delta \in [0,1)$ 表示违约后的回收比例, $Q(u)$ 表示与当前合约价值有关的违约强度.
        该方程展示了金融定价模型中非线性项的一个典型来源: 市场风险因素或信用风险因素会使得原本线性的定价方程变为非线性.

        \item
        在统一框架下, Fisher--KPP 方程可取
        \begin{equation*}
          \mu(t, x) = 0,\qquad
          \sigma(t, x) = \sqrt{2} I_{d},
        \end{equation*}
        以及
        \begin{equation*}
          f(t, x, y, z) = y(1 - y).
        \end{equation*}
        代入统一形式后得到终端值问题
        \begin{equation*}
          \frac{\partial u}{\partial t}(t, x)
          + \Delta u(t, x)
          + u(t, x)(1 - u(t, x))
          =0.
        \end{equation*}
        Fisher--KPP 方程是另一类经典的非线性反应--扩散方程, 最初用于描述种群扩散和有利基因传播等问题.
        其典型特点是反应项具有 Logistic 增长结构: 当解较小时近似线性增长, 当解接近饱和值时增长受到抑制.
        因此, 该方程常被用来刻画带有扩散机制的传播前沿问题.
        对应的正向时间形式为
        \begin{equation*}
          \frac{\partial v}{\partial s}(s, x) = \Delta v(s, x) + v(s, x)(1 - v(s, x)).
        \end{equation*}
        通过同样的时间反向变换 $u(t, x) = v(T - t, x)$, 可以得到上面的终端值问题.
        与 Allen--Cahn 方程类似, Fisher--KPP 方程也是反应--扩散型方程,
        但其非线性项对应的是 Logistic 型增长, 而不是双势阱结构.
      \end{enumerate}

  \section{随机特征方法}\label{sec:rfm}
    通常的神经网络通过多层线性变换与非线性激活函数的叠加来逼近目标函数, 这会导致需要求解一个非线性的非凸优化问题.
    随机特征法的核心思想是将内层参数随机初始化并固定, 仅对输出层参数进行优化,
    从而将原问题转化为一个凸优化问题, 在平方损失函数下进一步退化为最小二乘问题.
    这种方法通过减少需要优化的参数自由度来显著简化训练过程, 以牺牲一定的表达能力为代价换取更高的计算效率与稳定性.
    由于内层参数固定, 随机特征法本质上可以视为对核方法的一种近似,
    即通过随机特征映射将核函数表示为特征函数的数学期望, 并利用有限样本平均进行近似.

    在实际应用中, 随机特征的分布及尺度参数会显著影响逼近能力.
    通过合理选择这些参数可以使特征函数张成更丰富的函数空间, 从而提高数值求解的精度.
    因此, 其在保持较低优化复杂度的同时, 为高维函数逼近提供了一种有效的替代方案.
    特别地, 由于其线性结构, 随机特征法在基于最小二乘的偏微分方程数值方法中具有天然优势.

    除高维问题外, 随机特征方法也可以用于常规低维 PDE 的数值近似.
    在低维情形中, FDM, FEM 和 SM 等方法已经形成了较成熟的理论和实现框架.
    这些方法通常先对空间区域进行网格剖分, 或选取一组确定性的局部基函数,
    再结合时间推进格式或变分形式得到有限维代数系统.
    与此不同, 随机特征方法直接选取一组由随机参数确定的全局特征函数,
    并将解函数近似为
    \begin{equation*}
      u_H(x) = \sum_{j=1}^{H} \alpha_j \phi_j(x),
    \end{equation*}
    其中 $\phi_j$ 表示固定的随机特征函数, $\alpha_j$ 为待求的线性系数.
    对稳态椭圆方程, 可以在区域内部采样配置点并要求 PDE 残差尽可能小,
    同时在边界采样点上加入边界条件残差.
    对抛物型方程, 可以在时间和空间上共同采样,
    也可以在给定时间层上分别构造随机特征近似.
    当方程和边界条件关于 $u_H$ 是线性的时,
    最终得到的参数求解问题通常就是线性最小二乘问题或带正则化的线性最小二乘问题.

    因此, 随机特征方法可以看作一种无网格 (mesh-free) 的残差最小化方法.
    它不要求规则网格, 形式上也较容易推广到高维输入.
    但其数值效果依赖于特征函数的选取, 特征宽度, 采样点数量以及线性系统的条件数.
    在低维光滑问题上, 传统网格方法往往仍然更加稳定和高效;
    随机特征方法的作用更多体现在提供一种统一的函数逼近方式,
    尤其适合与随机采样和最小二乘求解结合.
    在本文采用的 BSDE 离散框架中, 随机特征并不直接近似 $u(t,x)$,
    而是用于近似每个时间层上的梯度相关变量 $Z_k$.
    这样既保留了 BSDE 路径采样避免高维网格的特点,
    又利用固定随机特征带来的线性输出层结构,
    使线性方程可以直接归结为最小二乘系统,
    半线性方程也只需在输出层参数上进行非线性迭代.

    \subsection{对比方法: Deep BSDE 方法}\label{subsec:deepbsde_intro}
      深度倒向随机微分方程方法 (Deep Backward Stochastic Differential Equation method, Deep BSDE method)
      是本文数值实验中采用的对比方法.
      Han, Jentzen 和 E\cite{hanSolvingHighdimensionalPartial2018} 从半线性抛物型 PDE 的 BSDE 表示出发,
      沿正向随机过程生成样本路径, 并用神经网络近似每个离散时间层上的梯度项
      $\sigma^{\operatorname T}(t_k,X_{t_k})\nabla u(t_k,X_{t_k})$.
      初值 $Y_0$ 与各时间层神经网络参数共同作为待优化变量,
      通过最小化终端误差
      \begin{equation*}
        Y_N - g(X_N)
      \end{equation*}
      的平方损失来训练模型.

      该方法的特点是把 PDE 的离散求解写成沿随机路径传播的优化问题.
      其中, 正向过程 $X_t$ 由给定的漂移项 $\mu$ 和扩散项 $\sigma$ 生成,
      倒向过程中的 $Z_t$ 则由神经网络给出.
      由于整个算法只需要在样本路径上计算, 不需要在高维空间中构造网格,
      因而可以用于传统网格方法难以直接处理的高维问题.
      文献\cite{hanSolvingHighdimensionalPartial2018} 在非线性 Black--Scholes 方程, HJB 方程和 Allen--Cahn 方程上的实验表明,
      这种框架能够在较高维度下给出有效的数值近似.
      这与 E\cite{eMachineLearningComputational2020} 对机器学习和计算数学关系的讨论是一致的:
      对高维 PDE 而言, 困难的核心逐渐从网格离散转向高维函数表示和随机优化.

      与本文采用的随机特征方法相比, Deep BSDE 方法需要优化的参数更多.
      在常见实现中, 每个时间层都需要用一个神经网络近似梯度函数,
      因此参数量会随时间剖分数和网络规模增加.
      训练过程需要对全部网络参数进行反向传播,
      对网络宽度, 学习率, 初始化和训练轮数等超参数也较为敏感.
      因此, 本文仅将 Deep BSDE 作为对比方法, 用于检验随机特征方法在相同 BSDE 离散框架下的计算效率与初值误差表现.

  \section{本文主要工作与章节安排}\label{sec:work}
    本文围绕 BSDE 离散框架下的随机特征方法展开, 目标是在不构造高维空间网格的前提下,
    用较低复杂度的参数求解过程近似高维线性与半线性抛物型方程的初值.
    主要工作包括以下几个方面.

    首先, 本文从统一的半线性抛物型方程出发,
    通过 Itô 公式将终值问题转化为对应的 BSDE 形式,
    并在离散时间层上构造样本路径.
    在此基础上, 采用随机特征函数近似 BSDE 中的梯度相关变量 $Z_k$,
    固定随机特征的内层参数, 只求解初值 $y_0$ 和输出层参数 $\alpha$.
    这种处理保留了路径采样适合高维问题的特点,
    同时避免了对多层神经网络全部参数进行训练.

    其次, 对生成元关于 $(y,z)$ 为线性的方程,
    本文推导了离散终端值关于 $y_0$ 和 $\alpha$ 的仿射表达式,
    将终端匹配问题化为带岭正则化的线性最小二乘问题.
    正则化只作用在随机特征输出层参数上, 不惩罚初值项,
    以减少病态线性系统对参数规模的影响.
    对半线性方程, 终端残差不再关于参数保持线性,
    本文将问题写成非线性最小二乘形式, 并采用 LM 方法迭代求解.
    为了提高迭代过程的可控性, 文中给出了终端残差关于参数的灵敏度递推,
    用于构造 LM 子问题中的雅可比矩阵.

    最后, 本文基于 C++ 和 libtorch 实现了上述求解流程,
    并在线性 Black--Scholes 方程, 热方程以及半线性 HJB-LQ 方程, Allen--Cahn 方程上进行数值实验.
    实验中同时报告训练路径上的终端残差和独立测试路径上的终端残差,
    并考察样本数, 随机特征宽度, 维度和随机种子对结果的影响.
    此外, Deep BSDE 方法仅作为对比方法出现在数值实验中,
    用于观察本文随机特征方法在相同 BSDE 路径采样框架下的计算时间和初值误差表现.

    后文结构安排如下.
    第\ref{ch:methodology}章给出本文方法的具体构造,
    包括 PDE 到 BSDE 的转化, 样本路径与随机特征的构造,
    线性方程的最小二乘求解, 半线性方程的 LM 求解以及测试误差评估.
    第\ref{ch:numerical}章给出四个高维算例的实验设置和数值结果,
    并分析样本数, 特征宽度, 维度变化和随机种子稳定性等因素.
    附录部分补充关键公式推导, 实验伪代码和参考值计算方法.
